\section{Threats to validity}
\label{sec:threats}
% Former subsection labels, kept so that existing cross-references resolve to this section; their
% content is now a row of Table~\ref{tab:threats}.
\label{subsec:rule5}\label{subsec:loop}\label{subsec:weights}\label{subsec:benchmarks}
\label{subsec:ourablation}\label{subsec:searchbias}\label{subsec:searchfreeze}\label{subsec:lessons}

Table~\ref{tab:threats} gives each threat's direction, the mitigation performed, and what
that mitigation leaves unestablished; several are weaker than we would have chosen. Three threats bound
the conclusions and are discussed below.

\begin{table}[tp]
\caption{Threats to validity. \emph{Direction} is the effect on the estimate named in the cell.
\emph{Mitigation} lists only checks that were performed. Row groups: review process; sample and search;
evidence base; taxonomy.}
\label{tab:threats}
\footnotesize
\begin{tabular}{@{}p{0.17\linewidth}p{0.10\linewidth}p{0.38\linewidth}p{0.27\linewidth}@{}}
\toprule
Threat & Direction & Mitigation performed & Does not establish \\
\midrule
\raggedright No human screener or coder (amendments 4, 5) & \raggedright inflates census size &
\raggedright Every coded value carries a verbatim-verified quote; 58,426 of 60,881 screening quotes (96.0\%)
verbatim; recall 27/27 on a reference set; independent pipeline on 434 shared records, agreement 0.843,
$\kappa = 0.664$; second model reading on 3,075 escalated records, $\kappa = 0.473$; tier~2 overturned 670 of 1,532
tier-1 excludes (43.7\%) and 140 of 1,543 sampled tier-1 includes (9.1\%) &
\raggedright That a human expert would decide the same way \tabularnewline
\addlinespace[2pt]
\raggedright Reliability is reproducibility & \raggedright unknown &
\raggedright 247 double-coded systems: agreement 0.870, mean per-dimension $\kappa$ 0.784, AC1 0.855; 37 of 38
dimensions at the 0.6 floor &
\raggedright Correctness: a shared systematic error is invisible to $\kappa$ \tabularnewline
\addlinespace[2pt]
\raggedright Coding rule rewritten mid-project (amendment 9) & \raggedright inflates reliability &
\raggedright Rewritten, validated on 217 systems, then applied; \texttt{not\_reported} 26.9\%$\to$48.3\%, agreement
0.597$\to$0.872, stacked $\kappa$ 0.564$\to$0.832; superseded coding released &
\raggedright That the rule was not tuned to the dimensions seen failing \tabularnewline
\addlinespace[2pt]
\raggedright \texttt{loop\_primitives} below the floor & \raggedright deflates associations &
\raggedright Kept and flagged: $\kappa$ 0.5904 [0.523, 0.657], AC1 0.6283, per-value $\kappa$ 0.731; cluster-bootstrap
intervals for all 38 dimensions &
\raggedright Which side of the floor four dimensions lie on; their intervals include 0.6 \tabularnewline
\midrule
\raggedright Sampling and weighting & \raggedright unknown; misses likely understate O &
\raggedright Design weights (P: 150 of 683 at 4.55; O: 100 of 4,837 at 48.37); standard errors up to 3.7
points on silence rates and 18.6 on value shares; weighted and unweighted rates reported &
\raggedright Frame completeness: an independent pipeline's includes were 308 of 411 (75\%) in our pool; of
the 103 outside it, 37 passed our full-text screen \tabularnewline
\addlinespace[2pt]
\raggedright English-language and venue bias & \raggedright deflates census size &
\raggedright Eight sources; the 102 independent includes absent from the whole harvest split into 36
Boolean-matchable and 66 vocabulary misses &
\raggedright How many harnesses are documented only in other languages \tabularnewline
\addlinespace[2pt]
\raggedright Search frozen at 2026-09-16 & \raggedright unknown &
\raggedright Each system pinned to a commit (amendment 8); dates released &
\raggedright Trends beyond cohort comparisons; 2026 is the modal release year (52\% weighted, 412 dated systems) \tabularnewline
\midrule
\raggedright Benchmark caveats & \raggedright deflates cross-paper effect &
\raggedright Within-study design cancels versioning, contamination, entanglement, and drift; cross-paper keys fix
benchmark, split, metric, and model &
\raggedright Drift attenuates the $+0.88$ within-key sd and baseline selection inflates it (\rqXPOtherEffect{} on the
\rqXPOtherKeys{} keys without recurring baselines), so its net bias direction is unknown; it is on another scale than
the within-study upper bounds \tabularnewline
\addlinespace[2pt]
\raggedright Author self-reporting & \raggedright inflates ablation effects &
\raggedright Sign distribution: \rqSignSharePct\% of pooled contrasts favour the component (67 of 75 on the
coded set); discount per dimension, median \rqMedianDiscountPct\%, which no longer separates dimensions &
\raggedright A correction; no external estimate of the bias exists \tabularnewline
\addlinespace[2pt]
\raggedright Our own filed ablation & \raggedright inflates in our favour &
\raggedright Protocol frozen before the first scored instance; a compute-matched arm that can return a null;
per-instance logs to be released &
\raggedright Any constraint tested by a result; only pilot data are collected \tabularnewline
\midrule
\raggedright Construct validity of the taxonomy & \raggedright unknown &
\raggedright Dimensions frozen before coding and crosswalked to prior taxonomies; one pair above $V = 0.6$
(0.720, $n = 841$) &
\raggedright That 38 dimensions carve the object; no loop-driver axis \tabularnewline
\bottomrule
\end{tabular}
\end{table}

\paragraph{No human screener or coder.}
\label{subsec:nohuman}
\label{subsec:repro}
As disclosed in \S\ref{subsec:llmdisclosure}, reliability is reproducibility, never correctness: two model readings can agree and both be
wrong in the same way. What stands in for correctness, the first row of Table~\ref{tab:threats}, works at
three levels: whether a system exists, whether it is in scope, and whether its quoted evidence is real.
None tests whether a coded value is the one an expert would choose. The second model reading
goes against us and makes the census an upper estimate. It covers an escalated, non-random subset:
the disputed part of the screen, not the whole.

\paragraph{Author self-reporting bounds RQ3 from above.}
\label{subsec:selfreport}
Every ablation contrast is an author measuring their own component, and the sign distribution shows why
that matters: a component that helped only sometimes would draw more dissent than one contrast in ten.
The discount therefore turns each pooled effect into an upper bound, not an estimate: ``at most about
\rqSVDiscounted{} relative for self-verification'', never ``\rqSVDiscounted''. Because most poolable
dimensions carry a bound of similar size, the bounds do not rank components. Once its chance baseline excluded the papers'
own ablation arms, the own-arm test stopped adding a second measure, so the bound rests on the sign
distribution alone. A correction needs an external estimate, such as an independent re-run of a random
sample of published baseline tables. Our own filed ablation is exposed to the same bias.

\paragraph{Construct validity of the taxonomy.}
\label{subsec:construct}
Thirty-eight dimensions are a choice. Our schema has no counterpart for the loop driver, the axis
\citet{rombaut2026insidescaffold} calls the most fundamental architectural distinction it draws, nor for
tool discovery strategy or multi-model routing. Choices coupled through an uncoded axis would produce the
same weak coupling we measure, a mean corrected Cram\'er's V of 0.167, so ``no families'' and ``wrong
axes'' are observationally equivalent on our cells. Separating them needs a second taxonomy applied to the
same systems. The released cell-level dataset permits one; extending ours would mean recoding 1,256
systems against a schema frozen before coding began.
